\subsection{Búsqueda Exponencial}\label{sec:exponencial}
La búsqueda exponencial es un algoritmo híbrido de búsqueda que combina una etapa inicial de expansión rápida del rango de búsqueda con una etapa posterior de búsqueda binaria. Su principal utilidad aparece cuando trabaja con arreglos ordenados y se desea localizar eficientemente un elemento sin conocer de antemano una zona acotada en la que este pueda encontrarse.

La idea central del algoritmo consiste en aumentar exponencialmente el índice inspeccionado hasta encontrar un valor mayor o igual al objetivo, o bien hasta superar el tamaño del arreglo. Una vez determinado ese intervalo, se aplica búsqueda binaria únicamente sobre dicho subarreglo. De esta manera, el algoritmo evita buscar desde el principio sobre todo el arreglo y acota rápidamente la región relevante.

En este proyecto, la búsqueda exponencial se utiliza sobre arreglos ordenados por \texttt{id}, de modo que la comparación entre elementos mantenga la propiedad necesaria para descartar regiones completas del arreglo.

\subsubsection{Funcionamiento general}
El procedimiento puede describirse en dos etapas:
\begin{enumerate}
    \item Expansión exponencial del rango: Se comienza revisando el primer elemento del arreglo. Si no coincide con el valor buscado, se inspeccionan posiciones que crecen como potencias de dos: \texttt{1, 2, 4, 8, 16, ...} hasta encontrar una posición cuyo valor sea mayor o igual al objetivo, o hasta alcanzar el final del arreglo.
    \item Búsqueda binaria en el intervalo encontrado: Una vez acotado el rango donde potencialmente se encuentra el elemento, se ejecuta una búsqueda binaria clásica sobre el subarreglo.
\end{enumerate}
Este enfoque resulta eficiente porque la fase exponencial reduce rápidamente el espacio de búsqueda, mientras que la fase binaria aprovecha esa cota para completar la localización del elemento.

\begin{algorithm}
    \caption{Búsqueda exponencial $\triangleright$ $V[0..n-1]$}
    \begin{algorithmic}[1]
        \Procedure{ExponentialSearch}{$V, n, x, comp\_f$}

        \If{$n \leq 0$}
            \State \Return $-1$
        \EndIf

        \If{$comp\_f(\&V[0], x) = 0$}
            \State \Return $0$
        \EndIf

        \State $bound \gets 1$

        \While{$bound < n$ \textbf{and} $comp\_f(\&V[bound], x) < 0$}
            \State $bound \gets bound \cdot 2$
        \EndWhile

        \State $beg \gets \dfrac{bound}{2}$

        \If{$bound < n$}
            \State $end \gets bound$
        \Else
            \State $end \gets n - 1$
        \EndIf

        \State \Return \Call{BinarySearch}{$V, beg, end, x, comp\_f$}

        \EndProcedure
    \end{algorithmic}
\end{algorithm}

La parte que destaca de este algoritmo es la expansión exponencial. En lugar de recorrer el arreglo elemento a elemento, el algoritmo duplica progresivamente la posición inspeccionada, lo que le permite acercarse muy rápido a la zona donde podría estar el valor buscado. Una vez encontrado el límite superior, el problema se reduce a una búsqueda binaria convencional en un intervalo acotado.

En este sentido, este algoritmo se puede interpretar como una optimización para situaciones en las que el valor buscado podría hallarse relativamente cerca del inicio del arreglo, pero sin perder eficiencia cuanto más lejos este el elemento.

\subsubsection{Análisis de complejidad}
La complejidad temporal de la búsqueda exponencial se obtiene sumando el costo de sus dos etapas, donde la primera corresponde a la fase exponencial, donde cada en cada iteración el índice se duplica, lo que corresponde a orden de $log$ $n$, y la segunda corresponde a la fase binaria que se ejecuta una vez determinado el intervalo, que también cuesta $log$ $n$.

Por lo tanto, el costo total es: $$O(log\text{ }n)$$
De esta forma, la complejidad temporal queda:
\begin{itemize}
    \item Mejor caso: $O(1)$, si el elemento se encuentra en primera posición.
    \item Caso promedio: $O(log\text{ }n)$
    \item Peor caso: $O(log\text{ }n)$
\end{itemize}

Aunque comparte el mismo orden de complejidad que búsqueda binaria, la búsqueda exponencial incorpora una fase previa adicional. Sin embargo, esa fase también crece logarítmicamente, por lo que el comportamiento asintótico no se ve alterado.

\subsubsection{Comentario general}
La búsqueda exponencial representa una extensión de la búsqueda binaria para escenarios en los que no se conoce de antemano una región acotada del arreglo donde podría encontrarse el elemento buscado. Su ventaja principal se encuentra en que acota rápidamente el intervalo antes de aplicar búsqueda binaria, manteniendo en todo momento una complejidad logarítmica. En el contexto de este proyecto, constituye una alternativa eficiente sobre arreglos ordenados por \texttt{id}, y resulta interesante para comparar cómo distintas variantes de búsqueda basadas en divide y vencerás alcanzan órdenes de crecimiento similares mediante estrategias operacionales distintas.